% Contenido del póster: Alvarez-Franco, Castañeda, Agudelo Zuluaga y Diaz (2026), Journal of Economics, Finance and
% Administrative Science. DOI 10.1108/jefas-09-2024-0280
% Elaborado con apoyo de IA a partir del resumen oficial del artículo (OpenAlex), 2026-11-12.
\def\DOI{10.1108/jefas-09-2024-0280}
\def\TituloEs{Determinantes de la inclusión financiera digital en América Latina y el Caribe}
\def\TituloOriginal{Determinants of digital financial inclusion in Latin America and the Caribbean}
\def\Autores{Alvarez-Franco, Castañeda, Agudelo Zuluaga y Diaz}
\def\Anio{2026}
\def\Revista{Journal of Economics, Finance and Administrative Science}

\def\Problema{Los servicios financieros digitales crecen en América Latina y el Caribe, pero su acceso y uso no se distribuyen por igual entre la población.}
\def\Objetivo{Investigar los determinantes de la inclusión financiera digital en la región, en particular la influencia de los factores sociodemográficos en el acceso y el uso de los servicios financieros digitales.}
\def\Metodo{Datos de la base Global Findex 2021 del Banco Mundial para 15 países de la región. Se construye un índice que agrupa las preguntas sobre acceso y uso de servicios financieros digitales y se estima un modelo probit de sus determinantes demográficos.}
\def\Resultados{\begin{itemize}
  \item La edad, la educación, el género, el ingreso, la situación laboral y el país de origen influyen de forma significativa en la inclusión financiera digital.
  \item Las mujeres y las personas con menos ingresos y educación están en desventaja.
  \item El estudio reconoce que no incluye otros determinantes, como la infraestructura, la regulación y los factores culturales.
\end{itemize}}
\def\Conclusion{Las políticas públicas deben reducir barreras como la alfabetización digital y el acceso a internet, y fortalecer la innovación fintech y la infraestructura en los países con menor adopción.}
\def\Relevancia{Las estadísticas de la SMV describen la oferta del financiamiento participativo, pero no quién lo usa. Este estudio muestra que la demanda de servicios financieros digitales depende de características de las personas; en el Perú, incorporar información de la demanda permitiría saber a quién llega el financiamiento participativo.}
\def\PalabrasClave{inclusión financiera digital, Global Findex, América Latina y el Caribe, modelo probit, fintech}
% --- Ficha del estudio (herramientas/posteres/ficha_articulos.py) ---
\def\Ambito{15 países de América Latina y el Caribe (Global Findex 2021)}
\def\Enfoque{Empírico: índice de inclusión y modelo probit}
\def\FiguraPeru{2}
% --- fin de la ficha ---
